\documentclass[12pt,a4paper]{report}
\usepackage[utf8]{inputenc}
\usepackage[brazilian]{babel}
\usepackage{geometry}
\usepackage{amsmath}
\usepackage{amssymb}
\usepackage{graphicx}
\usepackage{hyperref}
\usepackage{booktabs}
\usepackage{array}
\usepackage{float}
\usepackage{fancyhdr}\usepackage{helvet}
\usepackage{sectsty}
\usepackage{titlesec}
\usepackage{indentfirst}\usepackage{tabularx}
\renewcommand{\familydefault}{\sfdefault}
\allsectionsfont{\sffamily\fontsize{12}{14}\selectfont}

\titleformat{\chapter}[hang]{\normalfont\Large\bfseries}{\thechapter.}{1em}{}
\geometry{margin=2.5cm}

\setlength{\parskip}{0.5cm}
\setlength{\parindent}{0.75cm}

\pagestyle{fancy}
\fancyhf{}
\rhead{Relatório de Reclamações - Pombal}
\lhead{\thepage}
\renewcommand{\headrulewidth}{0.4pt}

\title{\textbf{RELATÓRIO DE RECLAMAÇÕES}\\
Município de Pombal - PB\\
Questões com a Concessionária de Energia}
\author{}
\date{\today}

\begin{document}

\maketitle

\tableofcontents
\newpage

\chapter{Objetivo}

O presente relatório tem por finalidade apresentar um panorama atualizado das reclamações administrativas em tramitação relacionadas às unidades consumidoras vinculadas ao município de \textbf{Pombal-PB}, indicando o assunto tratado, o objeto da demanda, a unidade consumidora envolvida e o atual status de análise. 

As reclamações encontram-se em acompanhamento perante a Ouvidoria da ANEEL e a Ouvidoria da concessionária, tratando principalmente de descumprimento de cobranças indevidas e pedidos de devolução complementar. 

\vspace{3cm}

\chapter{Resumo Geral das Reclamações}

Atualmente, foram identificadas 22 reclamações, a maioria concluída e algumas ainda em análise, como mostra a tabela \ref{tabela_resumo}.

\begin{table}[H]
\centering
\caption{Resumo de Reclamações por Situação}
\renewcommand{\arraystretch}{1.5}
\begin{tabularx}{16cm}{|>{\centering\arraybackslash}X|>{\centering\arraybackslash}c|}
\hline
\textbf{Situação} & \textbf{Quantidade} \\
\hline
Concluída & 15 \\
\hline
Em análise na ouvidoria da ANEEL & 1 \\
\hline
Pendente contestação & 6 \\
\hline
\textbf{Total} & \textbf{22} \\
\hline
\end{tabularx}
\label{tabela_resumo}
\end{table}

A maior parte das demandas ainda em tramitação estão com pendencia de contestação na ouvidoria da concessionária. 

\vspace{3cm}

\chapter{Principais Teses em Discussão}

\begin{table}[H]
\centering
\caption{Resumo de teses em discussão}
\renewcommand{\arraystretch}{1.5}
\begin{tabularx}{16cm}{|>{\centering\arraybackslash}X|>{\centering\arraybackslash}c|}
\hline
\textbf{Tese/assunto} & \textbf{Quantidade} \\
\hline
Enquadramento IP & 6 \\
\hline
Enquadramento AES & 12 \\
\hline
Não isenção de imposto de renda & 1 \\
\hline
Parcela DIC/2 & 1 \\
\hline
Valor incompativel com QIP & 1 \\
\hline
Perda nos reatores & 1 \\
\hline
\textbf{Total} & \textbf{22} \\
\hline
\end{tabularx}
\label{tabela_teses}
\end{table}

\vspace{3cm}

\chapter{Planilha Detalhada de Reclamações}

\section*{}

\begin{table}[H]
\centering
\begin{tabular}{|p{2cm}|p{3cm}|p{3cm}|p{3cm}|p{2cm}|}
\hline
\textbf{ID} & \textbf{Data} & \textbf{Descrição} & \textbf{Assunto} & \textbf{Status} \\
\hline
 &  &  &  &  \\
\hline
\end{tabular}
\end{table}

\vspace{3cm}

\chapter{Análise Geral}

\section*{}

\vspace{3cm}

\end{document}
